\rtmtight
\begin{tblr}{width=\textwidth,colspec={|Q[c,m,2.1cm]|X[l,m]|},hlines,vlines,hspan=minimal,rowsep=1.5pt,colsep=3pt,row{1}={bg=headgray,font=\bfseries}}
\SetCell{c,m}\hfil\logo\hfil & \textbf{POLITEKNIK NEGERI MALANG}\par\textbf{JURUSAN TEKNOLOGI INFORMASI}\par\textbf{PROGRAM STUDI: D4 TEKNIK INFORMATIKA} \\
\SetCell[c=2]{l} \textbf{RENCANA TUGAS MAHASISWA (RTM) DAN RUBRIK PENILAIAN --- DRAF KONSEP A: SOLUSI TI} & \\
Nama Mata Kuliah & Kuliah Kerja Nyata Tematik \\
Kode Mata Kuliah & RTI257002 (sementara) \quad \textbf{sks / Jam perminggu:} 2 SKS / 4 jam/minggu \quad \textbf{Semester:} 6/7 (sementara) \\
Dosen Pengampu & Tim Pengajar Program Studi D4 TI \\
\end{tblr}
\assessmentblock{Proposal Program Kerja KKN}{Case Method dan penyusunan dokumen}{SCPMK0101-06001, SCPMK0301-06002, SCPMK0803-06003}{Mahasiswa menyusun proposal program kerja KKN berdasarkan observasi dan analisis kebutuhan mitra, memuat rencana kegiatan, jadwal, dan rancangan solusi TI yang akan dibangun.}{Melakukan observasi lapangan, menganalisis kebutuhan mitra secara partisipatif, menyusun proposal, dan mempresentasikannya dalam seminar proposal.}{Dokumen proposal program kerja KKN dan bahan presentasi seminar proposal.}{Ketepatan analisis kebutuhan dan kelayakan program & 5 & Analisis kebutuhan mitra akurat dan mendalam; program kerja yang diusulkan realistis, relevan, dan menjawab masalah nyata mitra. & Analisis kebutuhan cukup memadai; program kerja relevan tetapi kurang detail pada beberapa aspek pelaksanaan. & Analisis kebutuhan dangkal atau tidak berbasis data lapangan; program kerja tidak relevan atau tidak realistis. \\ \\
Kejelasan rancangan solusi TI & 5 & Rancangan solusi TI dijelaskan secara rinci, sesuai kebutuhan mitra, dan layak untuk diimplementasikan pada tahap pelaksanaan. & Rancangan solusi TI cukup jelas tetapi beberapa bagian teknis masih perlu disempurnakan. & Rancangan solusi TI tidak jelas, tidak sesuai kebutuhan mitra, atau tidak layak diimplementasikan. \\ \\
Kualitas presentasi dan etika penyusunan & 5 & Presentasi proposal terstruktur dan meyakinkan; etika penyusunan (kejujuran data, pengakuan sumber) terjaga penuh. & Presentasi cukup baik dengan sedikit kekurangan struktur; etika penyusunan secara umum terjaga. & Presentasi tidak terstruktur atau terdapat pelanggaran etika penyusunan proposal. \\}

\assessmentblock{Produk/Solusi TI untuk Mitra}{Portofolio dan praktik lapangan}{SCPMK0803-06003, SCPMK0607-06004}{Mahasiswa merancang, mengimplementasikan, dan menyerahterimakan solusi teknologi informasi kepada mitra sesuai kebutuhan yang telah disepakati dalam proposal program kerja.}{Mendesain solusi TI, mengimplementasikan dan mengujinya di lokasi mitra, melatih pengguna, serta menyempurnakan solusi berdasarkan umpan balik monev dan mitra.}{Solusi TI yang berjalan (aplikasi/sistem/purwarupa) beserta dokumentasi pengembangan dan berita acara serah terima ke mitra.}{Kesesuaian solusi dengan kebutuhan mitra & 15 & Solusi TI berfungsi penuh dan secara langsung menjawab kebutuhan mitra yang dirumuskan di proposal, disertai bukti penggunaan oleh mitra. & Solusi TI berfungsi dan menjawab sebagian besar kebutuhan mitra; beberapa fitur belum optimal atau belum digunakan mitra. & Solusi TI tidak berfungsi dengan baik atau tidak sesuai kebutuhan mitra yang telah disepakati. \\ \\
Kualitas rekayasa dan implementasi teknis & 15 & Implementasi rapi, teruji, dan terdokumentasi; menerapkan praktik rekayasa perangkat lunak yang baik (desain, pengujian, dokumentasi kode). & Implementasi berfungsi tetapi kualitas rekayasa (pengujian/dokumentasi kode) masih terbatas. & Implementasi tidak stabil, tidak teruji, atau tidak dapat dijelaskan secara teknis oleh mahasiswa. \\ \\
Pengelolaan pelaksanaan dan keterlibatan mitra & 10 & Pelaksanaan program terjadwal dengan baik, mitra terlibat aktif dalam pengembangan dan pelatihan penggunaan solusi. & Pelaksanaan cukup terjadwal; keterlibatan mitra ada tetapi tidak konsisten sepanjang program. & Pelaksanaan tidak terjadwal atau mitra tidak dilibatkan dalam proses pengembangan solusi. \\}

\assessmentblock{Laporan KKN}{Case Method dan penulisan ilmiah}{SCPMK0101-06001, SCPMK0803-06003, SCPMK0607-06004}{Mahasiswa menyusun laporan KKN yang komprehensif mencakup profil mitra, pelaksanaan program kerja, dokumentasi teknis solusi TI, dan refleksi etika serta tanggung jawab profesional.}{Mengumpulkan data dan dokumentasi selama pelaksanaan program, menyusun laporan sesuai sistematika Polinema, bimbingan dengan DPL, dan revisi berdasarkan masukan.}{Naskah laporan KKN lengkap sesuai pedoman Polinema, dilengkapi lampiran pendukung (logbook, dokumentasi teknis solusi, berita acara serah terima).}{Kelengkapan struktur dan dokumentasi pelaksanaan & 10 & Laporan memuat semua bab yang dipersyaratkan dengan dokumentasi pelaksanaan program kerja yang lengkap dan akurat. & Struktur laporan lengkap tetapi beberapa bagian dokumentasi pelaksanaan kurang mendalam. & Struktur laporan tidak lengkap atau dokumentasi pelaksanaan tidak mencerminkan kegiatan yang sesungguhnya. \\ \\
Kualitas dokumentasi teknis solusi TI & 8 & Dokumentasi teknis solusi TI (rancangan, implementasi, pengujian) lengkap, jelas, dan dapat digunakan untuk pemeliharaan lanjutan oleh mitra. & Dokumentasi teknis cukup lengkap tetapi ada bagian yang kurang detail. & Dokumentasi teknis solusi TI minim atau tidak dapat ditelusuri oleh pembaca. \\ \\
Refleksi etika dan tanggung jawab profesional & 7 & Refleksi menunjukkan pemahaman mendalam terhadap etika kerja lapangan dan tanggung jawab profesional dalam melayani kebutuhan mitra. & Refleksi cukup informatif tetapi kedalaman analisis etika masih terbatas. & Refleksi bersifat deskriptif semata tanpa analisis etika atau tanggung jawab profesional. \\}

\assessmentblock{Seminar Hasil KKN}{Presentasi dan defense}{SCPMK0301-06002, SCPMK0607-06004}{Mahasiswa mempresentasikan dan mempertahankan laporan KKN beserta solusi TI yang dihasilkan di hadapan tim penguji, menunjukkan penguasaan atas proses dan hasil program.}{Menyiapkan bahan presentasi, memaparkan ringkasan program kerja dan solusi TI, mendemonstrasikan solusi, serta menjawab pertanyaan penguji secara analitis.}{Presentasi sidang (slide dan demo solusi TI), naskah laporan akhir yang telah direvisi, dan berita acara seminar hasil.}{Kualitas presentasi dan demonstrasi solusi & 7 & Presentasi terstruktur dan demonstrasi solusi TI berjalan lancar, meyakinkan penguji atas kualitas hasil program. & Presentasi cukup terstruktur; demonstrasi solusi berjalan tetapi kurang lancar atau kurang meyakinkan. & Presentasi tidak terstruktur atau demonstrasi solusi gagal/tidak dapat ditampilkan. \\ \\
Kemampuan menjawab pertanyaan penguji & 7 & Semua pertanyaan dijawab dengan tepat, didukung argumen logis dan bukti dari pengalaman lapangan serta solusi yang dibangun. & Sebagian besar pertanyaan dijawab dengan cukup baik, tetapi beberapa jawaban kurang didukung bukti. & Banyak pertanyaan tidak dapat dijawab atau tidak relevan dengan laporan dan solusi TI. \\ \\
Keterkaitan solusi TI dengan kebutuhan mitra & 6 & Mahasiswa dapat menjelaskan secara eksplisit keterkaitan solusi TI dengan kebutuhan dan dampak nyata bagi mitra. & Keterkaitan dijelaskan tetapi masih umum atau kurang spesifik terhadap dampak bagi mitra. & Tidak dapat menjelaskan keterkaitan solusi TI dengan kebutuhan mitra, atau pembahasan sangat dangkal. \\}

\begin{tblr}{width=\textwidth,colspec={|Q[l,m,3.2cm]|X[l,m]|},hlines,vlines,hspan=minimal,row{1-3}={bg=headgray},rowsep=1.5pt,colsep=3pt}
Jadwal Pelaksanaan: & Minggu 1--16 sesuai RPS. Setiap evaluasi memiliki RTM dan rubrik multi-kriteria. \\
Lain-Lain Yang Diperlukan: & LMS, logbook digital, perangkat lapangan mitra, dan perangkat presentasi. \\
Pustaka: & Mengacu pustaka utama dan pendukung pada bagian RPS. \\
\end{tblr}
